%!TEX root = ../../main.tex
% ============================================================
% estrutura/pos-textuais/indice-remissivo.tex
%
% Índice remissivo (NBR 6034), último elemento pós-textual.
%
% \printindex lê main.ind (gerado pelo makeindex) e o insere
% dentro do ambiente `theindex`. A classe abntex2 (via memoir)
% já cuida de tudo sozinha dentro de `theindex`:
%   • título "ÍNDICE REMISSIVO" no estilo de capítulo ABNT
%     (\indexname traduzido em configuracoes/utfpr-setup.tex)
%   • entrada automática no sumário
%   • layout em duas colunas (padrão de índice remissivo)
% Não envolva \printindex em multicols/twocolumn manualmente:
% o ambiente já faz \twocolumn internamente, e aninhar seria
% redundante (e provavelmente quebraria o layout de página).
%
% Fluxo de compilação necessário para o índice aparecer:
%   1. lualatex main        → gera main.idx (a partir dos \index{})
%   2. makeindex main.idx   → gera main.ind (lista ordenada)
%   3. lualatex main        → lê main.ind e imprime o índice
%   4. lualatex main        → resolve sumário/referências cruzadas
% No CI (.github/workflows/latex-release.yml) isso é automático:
% latexmk (latexmk_use_lualatex: true) detecta o .idx e chama o
% makeindex sozinho a cada build. Localmente, use `latexmk -lualatex
% main.tex` pelo mesmo motivo — rodar `lualatex main.tex` sozinho
% NÃO atualiza o índice.
% ============================================================

\printindex
